\section{Introduction}


% Critiques:
% focus on growth, not claiming identity (greatest accom, i am CS)
% dont undercut your published work
% second paragraph is too dense (keep high-level)
% -> what exactly is the problem
% stronger throughline: why these projects
% show, don't tell

% TODO: Additionally, the student will add a single overall introduction (300-words or more) that motivates
% why they selected those three projects as exemplary of their accomplishments in the program.


% Graduate school is research training. My published work is evidence that the training took, but partial evidence, since it was done with an advisor in areas where I already had relevant coursework. The projects in this portfolio are fuller evidence. They are my most technically demanding work, they were done without guidance, and together they span the full range of what I think computer science research requires: careful theory, efficient implementation, and the experimental loop between them. I selected these three because no single project demonstrates all of that; together they do.
% Each addresses a distinct problem in performance engineering of hash-based data structures and hash functions. The first explores whether hybridizing binary fuse filters can yield theoretically and practically relevant space-time tradeoffs. The second introduces an optimal parameter tuning algorithm for secure learned Bloom filters, minimizing expected false positive rate subject to a memory budget, resolving the major open direction left by the original paper. The third extends entropy-aware hash function specialization to new cost models and more efficient implementations. Each required building fluency in an area I had not worked in before, without coursework or an advisor to shortcut that process.
% The range across these projects is intentional. In the binary fuse filter work, large-scale empirical searches revealed that a linear blend of three and four hash functions traced a new Pareto frontier, a result that led directly to a conjecture in graph theory. The Bloom filter work required developing and solving an optimization problem from scratch. The hash function work required both theoretical analysis and low-level implementation. Each project also left open problems I continue to pursue, which is perhaps the best evidence that the work was real.
% They are the clearest evidence that I have become a computer scientist.

The projects in this portfolio are my greatest accomplishments from my graduate studies. Unlike my published research, which developed within an established 
advising relationship and with relevant coursework as a foundation, this work grew out of entirely self-directed study. I acquired the necessary background, 
conducted literature reviews, then defined and pursued research questions independently. These three projects are the culmination of that process. \\

\noindent
Each addresses a distinct problem in performance engineering of hash-based data structures and hash functions. The first explores whether hybridizing binary 
fuse filters can yield theoretically and practically relevant space-time tradeoffs. The second introduces an optimal parameter tuning algorithm for secure 
learned Bloom filters, minimizing expected false positive rate subject to a maximum false positive rate and memory budget. The third extends entropy-aware hash 
function specialization to new cost models and more efficient implementations. \\

\noindent
Together these projects move between theory and efficient implementation, driven by an experimental approach: data suggested new theory, and that theory 
motivated new experiments. Each required building fluency in an unfamiliar area, and each left open problems I continue to pursue. They are the clearest 
evidence that I have become a computer scientist.


% I believe my projects in performance engineering of hash function and hash-based are my greatest accomplishment during my graduate studies. Unlike my published 
% research, I had no grounding in coursework or advising. Instead, I independently acquired enough domain knowledge to state and pursue research questions. These 
% three projects are the culmination of my self-directed studies. \\


% strongest overarching theme:
% extract and tune parameters to improve performance
% 
% why these projects:
% unified domain (hashing and hash-based data structures)
% 3 elible projects I enjoyed the most and am most proud of (although I have other work with greater research contribution)
% highlight a variety of skills
% combination of performance engineering, applied theory, and machine learning

% bigdata:
% skills: parallel computing, cluster computing, parameter tuning, cache locality optimization
% future: prove pareto frontier of non uniform filters

% mlsec
% skills: deep learning, dynamic programming, parameter tuning
% future: extend to account for other factors

% entropy aware hashing:
% skills: vector instructions, alignment, pipelining, performance engineering, hardness, approximation, 
% future: prove hardness + approximation

% TODO: Additionally, the student will add a single overall introduction (300-words or more) that motivates
% why they selected those three projects as exemplary of their accomplishments in the program.



% Graduate school is research training. My published work is evidence that the training took, but partial evidence, since it was done with an advisor in areas where I already had relevant coursework. The projects in this portfolio are fuller evidence. They are my most technically demanding work, they were done without guidance, and together they span the full range of what I think computer science research requires: careful theory, efficient implementation, and the experimental loop between them. I selected these three because no single project demonstrates all of that; together they do.
% Each addresses a distinct problem in performance engineering of hash-based data structures and hash functions. The first explores whether hybridizing binary fuse filters can yield theoretically and practically relevant space-time tradeoffs. The second introduces an optimal parameter tuning algorithm for secure learned Bloom filters, minimizing expected false positive rate subject to a memory budget, resolving the major open direction left by the original paper. The third extends entropy-aware hash function specialization to new cost models and more efficient implementations. Each required building fluency in an area I had not worked in before, without coursework or an advisor to shortcut that process.
% The range across these projects is intentional. In the binary fuse filter work, large-scale empirical searches revealed that a linear blend of three and four hash functions traced a new Pareto frontier, a result that led directly to a conjecture in graph theory. The Bloom filter work required developing and solving an optimization problem from scratch. The hash function work required both theoretical analysis and low-level implementation. Each project also left open problems I continue to pursue, which is perhaps the best evidence that the work was real.
% They are the clearest evidence that I have become a computer scientist.